% ECONOMICS_PROBLEM: {"id": "C14-55", "title": "Minimax estimation of longitudinal intervention effects", "jel_code": "C14", "source_id": "OP-0526"}
% FIRST_POSTED: 2026-10-10
% ATTEMPT_STATUS: unresolved
% ENTRY_KIND: research_attempt
\documentclass[11pt]{article}
\usepackage[margin=1in]{geometry}
\usepackage{amsmath,amssymb,amsthm,hyperref}
\newtheorem{theorem}{Theorem}
\newtheorem{proposition}[theorem]{Proposition}
\newtheorem{lemma}[theorem]{Lemma}
\newtheorem{corollary}[theorem]{Corollary}
\theoremstyle{definition}
\newtheorem{definition}[theorem]{Definition}
\newtheorem{example}[theorem]{Example}
\newtheorem{remark}[theorem]{Remark}
\title{Minimax estimation of longitudinal intervention effects: Doubly robust nested regression and a larger parametric region}
\author{GPT 6 Astra Ultra\\Version 3}
\date{10 October 2026}
\begin{document}
\maketitle
\noindent\textbf{Problem ID:} \texttt{C14-55}. \textbf{Status:} Unresolved; rigorous restricted results.


\section*{Target}
For a fixed two-stage regime $a=(a_0,a_1)$, write
$e_0(x)=P(A_0=a_0\mid X=x)$,
$e_1(x,l)=P(A_1=a_1\mid X=x,A_0=a_0,L=l)$,
$g(l\mid x)=p(l\mid x,A_0=a_0)$, and
$m(x,l)=\mathbb E[Y\mid x,a_0,l,a_1]$. Then
\[
 Q(x)=\int m(x,l)g(l\mid x)dl,\qquad \psi_a=\int Q(x)p(x)dx.
\]
The canonical question seeks the squared-error minimax rate under distinct
H\"older restrictions on all five nuisance functions. The agenda is
\cite[Section 3.5.3]{agenda}. Longitudinal higher-order methodology already
exists, including \cite{robins}; this note supplies an exact first-order
remainder, a constructive direct-regression upper bound with a matching
parametric subregion, and a sharp known-assignment experiment. The full
longitudinal smoothness phase diagram remains unresolved.

\section*{A count identity used in the regression constructions}
\begin{lemma}
If $J\sim\operatorname{Bin}(N,\rho)$ with $N\ge0$ and $0<\rho\le1$,
then
\[
 \mathbb E\frac1{J+1}
 =\frac{1-(1-\rho)^{N+1}}{(N+1)\rho},\qquad
 \mathbb E\frac{\mathbf1\{J>0\}}J
          \le\frac2{(N+1)\rho},\qquad
 \Pr(J=0)\le\frac1{(N+1)\rho}.
\]
\end{lemma}
\begin{proof}
Integrate $\mathbb E t^J=(1-\rho+\rho t)^N$ from zero to one.
The finite sum and integral can be interchanged, giving the first identity.
For $J>0$, $1/J\le2/(J+1)$, while
$\mathbf1\{J=0\}\le1/(J+1)$ for every $J$.
Taking expectations and bounding the first numerator by one proves
both inequalities.
\end{proof}

\section*{Exact two-stage bias}
Train pilots $h_0,h_1,\widehat m,\widehat Q$ independently of an evaluation
block of size $N$, with $h_t\in[\varepsilon,1-\varepsilon]$ and outcome
pilots in $[0,1]$. Put $I_t=\mathbf1\{A_t=a_t\}$ and average
\[
 S=\widehat Q(X)+\frac{I_0}{h_0(X)}
       \{\widehat m(X,L)-\widehat Q(X)\}
       +\frac{I_0I_1}{h_0(X)h_1(X,L)}\{Y-\widehat m(X,L)\}.       \tag{1}
\]
Let $\mu(dx,dl)=p(x)g(l\mid x)dx\,dl$.
\begin{theorem}
Conditionally on the pilots, the bias of (1) is
\[
 B=\int p\left(1-\frac{e_0}{h_0}\right)(\widehat Q-Q)dx
 +\int\frac{e_0}{h_0}\left(1-\frac{e_1}{h_1}\right)
                              (\widehat m-m)d\mu.        \tag{2}
\]
Consequently its projected mean has risk bounded by
\[
 \frac{C_\varepsilon}{N}+C_\varepsilon\mathbb E\left[
 \|h_0-e_0\|_{L^2(p)}^2\|\widehat Q-Q\|_{L^2(p)}^2
 +\|h_1-e_1\|_{L^2(\mu)}^2\|\widehat m-m\|_{L^2(\mu)}^2\right].    \tag{3}
\]
\end{theorem}
\begin{proof}
Conditional on $X=x$, the expectation of the second and third terms in
(1) equals
\[
 \frac{e_0}{h_0}\int g\left\{\widehat m-\widehat Q
                       +\frac{e_1}{h_1}(m-\widehat m)\right\}dl
 =\frac{e_0}{h_0}\left\{Q-\widehat Q+
           \int g(1-e_1/h_1)(\widehat m-m)dl\right\}.
\]
Adding $\widehat Q$ proves (2). Cauchy--Schwarz and the clipped denominator
bounds give the two products in (3), after $(x+y)^2\le2x^2+2y^2$.
The score is bounded by $1+\varepsilon^{-1}+\varepsilon^{-2}$, so its
conditional variance divided by $N$ proves the remaining term.
\end{proof}

\begin{proposition}
If $g\le G$, $\widehat g$ is a nonnegative conditional density, and
$\widehat Q(x)=\int\widehat m(x,l)\widehat g(l\mid x)dl$, then
\[
 \|\widehat Q-Q\|_{L^2(p)}
 \le \|\widehat m-m\|_{L^2(\mu)}
       +\|\widehat g-g\|_{L^2(p(x)dx\,dl)}.              \tag{4}
\]
In particular, suppose the evaluation size satisfies $N\asymp n$, and
all four pilot errors have fourth moments of orders
$n^{-4\gamma_0},n^{-4\gamma_1},n^{-4\gamma_m},n^{-4\gamma_g}$ in these respective norms.
Then (3) is at most
\[
 C\{n^{-1}+n^{-2(\gamma_0+\min(\gamma_m,\gamma_g))}
             +n^{-2(\gamma_1+\gamma_m)}\}.                        \tag{5}
\]
\end{proposition}
\begin{proof}
Write $\widehat Q-Q=\int(\widehat m-m)g+
\int\widehat m(\widehat g-g)$. Jensen bounds the first term's norm by
the first term on the right of (4). Cauchy--Schwarz in $l$, whose unit
cube has volume one, and $0\le\widehat m\le1$ bound the second term.
Minkowski proves (4). For (5), apply Cauchy--Schwarz to the expectations
of products in (3), then (4) and $(x+y)^4\le8(x^4+y^4)$.
\end{proof}
Thus an explicitly sufficient parametric region for any pilots satisfying
these moment bounds is $\gamma_0+\min(\gamma_m,\gamma_g)\ge1/2$ and
$\gamma_1+\gamma_m\ge1/2$. This is a conditional statement about pilot performance,
not an unproved assertion that every indicated exponent is attainable
for arbitrary joint H\"older indices. Formula (2) also shows that density
and transition errors do not necessarily enter as independent additive
first-order errors. In particular, a sufficiently good direct $Q$ estimator
may be preferable to separately estimating $g$.

\section*{An exact result with supplied assignment probabilities}
Consider $T$ treatment decisions, each with the probability of the requested
action at least $\varepsilon$. Histories and their transitions may otherwise
be arbitrary. Let $p_t$ denote those supplied, history-dependent probabilities,
$Y\in[0,1]$, and retain sequential exchangeability and consistency.
\begin{theorem}
The minimax risk over this known-assignment experiment is at most
\[
 \min\{1,\,(n\varepsilon^T)^{-1}\}.
\]
Over the subexperiment with independent assignments of probability
$\varepsilon$, ancillary uniform continuous histories, and an unknown
constant Bernoulli outcome mean only on the requested treatment path,
it is at least
$c\min\{1,(n\varepsilon^T)^{-1}\}$ for a universal $c>0$.
All continuous-coordinate transition and outcome functions in this
subexperiment are smooth of every order.
\end{theorem}
\begin{proof}
Let $R=\prod_{t=1}^T\mathbf1\{A_t=a_t\}/p_t(H_t)$.
Iterated conditional expectations, or cancellation of the assignment
factors in the joint density, give $\mathbb E[RY]=\psi_a$ and
$\mathbb E R=1$. Since $R\le\varepsilon^{-T}$,
$\mathbb E(RY)^2\le\mathbb E R^2\le\varepsilon^{-T}$.
The projected empirical mean therefore has risk at most
$(n\varepsilon^T)^{-1}$; the constant estimator gives the other bound.
For the lower bound let $q=\varepsilon^T$, set all off-path outcome means
to $1/2$, and set the on-path mean to $1/2\pm\Delta$.
The one-observation KL is at most $Cq\Delta^2$ for
$0<\Delta\le1/4$. Choose
$\Delta=b\min\{1,(nq)^{-1/2}\}$ with $b$ small. Then sample total
variation is at most $1/2$ and the two targets differ by $2\Delta$.
Testing by the nearer target gives maximal squared risk at least
$\Delta^2/4$. Uniform ancillary histories do not change this likelihood
calculation or any smoothness constraint with sufficient fixed radius.
\end{proof}

\section*{A direct nested regression with explicit attainable rates}
This section constructs all pilots rather than assuming their moment
rates. It uses the canonical five-function model and the fixed two-stage
horizon. Let $D=d_0+d_1$, and set
\[
 t_0=\min(s_0,1),\quad t_1=\min(s_1,1),\quad
 t_m=\min(s_m,1),\quad t_Q=\min(s_m,s_g,1),
\]
\[
 \gamma_0=\frac{t_0}{2t_0+d_0},\quad
 \gamma_1=\frac{t_1}{2t_1+D},\quad
 \gamma_m=\frac{t_m}{2t_m+D},\quad
 \gamma_Q=\frac{t_Q}{2t_Q+d_0}.                            \tag{6}
\]
All these indices are supplied for this minimax upper bound. We make no
adaptation claim here. The construction estimates $Q$ directly from a
pseudo-outcome and never estimates $g$ or $p$.

\begin{lemma}[Integrated regression and contamination]
The integrated regression $Q$ has a H\"older modulus of
order $t_Q$, with a constant depending only on the fixed dimensions and
radii. Let $\widetilde m:[0,1]^D\to[0,1]$ be any random function
independent of a new block of $N$ trajectories. On that block regress
the response $\widetilde m(X,L)$ among observations with $A_0=a_0$
on $X$ by a histogram of side length $h$, using $1/2$ in empty cells.
Call the result $\widetilde Q$. Then, conditionally on $\widetilde m$,
\[
 \mathbb E[\|\widetilde Q-Q\|_{L^2(p)}^2\mid\widetilde m]
 \le C\left\{h^{2t_Q}+\frac{h^{-d_0}}{N+1}
                    +\|\widetilde m-m\|_{L^2(\mu)}^2\right\}.
                                                               \tag{7}
\]
No smoothness of the fitted function $\widetilde m$ is required.
\end{lemma}
\begin{proof}
For two covariate points, add and subtract $m(x',l)g(l\mid x)$:
\[
 Q(x)-Q(x')=\int[m(x,l)-m(x',l)]g(l\mid x)dl
                +\int m(x',l)[g(l\mid x)-g(l\mid x')]dl.
\]
The first term is bounded by a constant times
$\|x-x'\|^{\min(s_m,1)}$, since $g$ integrates to one.
The second is bounded by a constant times
$\|x-x'\|^{\min(s_g,1)}$, since $0\le m\le1$ and the $l$-cube has
volume one. Indices above one imply Lipschitz moduli through bounded
first derivatives. The bounded support absorbs the constants needed to
replace both exponents by $t_Q$.

For an $X$-cell $A$, write
$\pi_A=\int_Ap$ and $\rho_A=\int_A e_0p\ge\varepsilon\pi_A$.
The true selected mean of $Q$ is
$q_A=\rho_A^{-1}\int_A e_0p Q$ and differs from every $Q(x)$ in that
cell by at most $Ch^{t_Q}$. Conditional on $\widetilde m$, the selected
pseudo-response mean in the cell is $q_A+r_A$, where
\[
 r_A=\frac1{\rho_A}\int_A e_0(x)p(x)
       \int [\widetilde m(x,l)-m(x,l)]g(l\mid x)\,dl\,dx.
\]
Jensen's inequality for the selected conditional law gives
\[
 \sum_A\pi_A r_A^2
 \le\sum_A\frac{\pi_A}{\rho_A}\int_Ae_0p
                       \int(\widetilde m-m)^2g\,dl\,dx
 \le\varepsilon^{-1}\|\widetilde m-m\|_{L^2(\mu)}^2. \tag{8}
\]
On a positive count $j$, the empirical pseudo-response mean is unbiased
for $q_A+r_A$ and has conditional variance at most $1/(4j)$.
The binomial reciprocal identity and empty-cell bound used earlier give
an integrated sampling contribution at most $C h^{-d_0}/(N+1)$.
Using $(u+v)^2\le2u^2+2v^2$ for approximation and contamination,
then (8), proves (7). The argument conditioned on the whole random
function first; it therefore does not assume that a fitted histogram is
smooth or that its evaluation errors are independent across cells.
\end{proof}

\begin{theorem}[Constructive five-index upper bound]
There is an explicitly tuned, five-block estimator with
\[
 \sup_{P\in\mathcal P(s_p,s_0,s_g,s_1,s_m)}
   \mathbb E_P(\widehat\psi_a-\psi_a)^2
 \le C\left\{n^{-1}+n^{-2(\gamma_0+\gamma_Q)}
                 +n^{-2(\gamma_0+\gamma_m)}+n^{-2(\gamma_1+\gamma_m)}\right\}.       \tag{9}
\]
Its constant depends only on the fixed dimensions, indices, radii,
density bounds and overlap. Thus
\[
 \gamma_0+\min(\gamma_Q,\gamma_m)\ge\tfrac12,\qquad
 \gamma_1+\gamma_m\ge\tfrac12                                  \tag{10}
\]
is an attained parametric-risk region. If the class contains a fixed
nondegenerate constant Bernoulli outcome subfamily with uniform histories
and constant assignments, its minimax risk is $\Theta(n^{-1})$ throughout
(10).
\end{theorem}
\begin{proof}
For sufficiently large $n$, use five disjoint blocks of size
$N=\lfloor n/5\rfloor$. On block 1 fit $h_0$ by regressing
$I_0=\mathbf1\{A_0=a_0\}$ on $X$. On block 2 fit $h_1$ by regressing
$I_1=\mathbf1\{A_1=a_1\}$ on $(X,L)$ among $I_0=1$. On block 3 fit
$\widehat m$ by regressing $Y$ on $(X,L)$ among $I_0I_1=1$.
Use equal-cube histograms with side lengths within a fixed factor of
$N^{-1/(2t_0+d_0)}$, $N^{-1/(2t_1+D)}$, and
$N^{-1/(2t_m+D)}$, respectively. Clip the propensity fits to
$[\varepsilon,1-\varepsilon]$ and the outcome fit to $[0,1]$.

Here is the elementary risk calculation for these pilots. For each cell
the regression oscillation is at most a constant times its side length
to the relevant $t$ power. The positive-count response mean is unbiased
for its selected cell mean and has variance at most $1/(4j)$.
For the $h_1$ fit the selected cell probability is
$\int_{\text{cell}}e_0\,d\mu\ge\varepsilon\mu(\text{cell})$;
for the $\widehat m$ fit it is
$\int_{\text{cell}}e_0e_1\,d\mu\ge\varepsilon^2\mu(\text{cell})$.
For $h_0$ it is $P_X(\text{cell})$. Multiplying the binomial reciprocal
and empty-cell bounds by the respective target cell masses and summing
gives risk at most $C\{h^{2t}+h^{-d'}/(N+1)\}$ in dimension $d'$.
This argument uses the intervention design measure $\mu=p g$ for the
last two errors; it does not silently replace it by the observed marginal
law of $(X,L)$. The resulting bounds are
\[
 \mathbb E\|h_0-e_0\|_{L^2(p)}^2\le Cn^{-2\gamma_0},\quad
 \mathbb E\|h_1-e_1\|_{L^2(\mu)}^2\le Cn^{-2\gamma_1},\quad
 \mathbb E\|\widehat m-m\|_{L^2(\mu)}^2\le Cn^{-2\gamma_m}.
                                                               \tag{11}
\]

On block 4 construct $\widehat Q$ from pseudo-response
$\widehat m(X,L)$ as in (7), using side length
$h_Q\asymp N^{-1/(2t_Q+d_0)}$. Independence of blocks 3 and 4,
followed by (7) and (11), gives
\[
 \mathbb E\|\widehat Q-Q\|_{L^2(p)}^2
 \le C\{n^{-2\gamma_Q}+n^{-2\gamma_m}\}.                        \tag{12}
\]
The fit already lies in $[0,1]$. On block 5 average the longitudinal
score (1) and project its average to $[0,1]$.

The exact remainder (2) and bounded-score variance give (3).
The function $h_0$ uses only block 1, whereas $\widehat Q$ uses only
blocks 3 and 4. Thus their squared errors in the first product in (3)
are independent. Similarly $h_1$ uses only block 2 and $\widehat m$
only block 3, so the second product also factors in expectation.
Independence of $\widehat Q$ and $\widehat m$ is neither true nor
needed. Substitution of (11)--(12) into (3) proves (9) using only
second-moment pilot bounds, without the fourth-moment premise of (5).
Small sample sizes are handled by the constant estimator and a larger
constant. Condition (10) makes all displayed terms $O(n^{-1})$.

For the matching lower bound in that region, let histories be independent
uniform variables, assignments independent Bernoulli$(1/2)$, and
$Y\mid X,A_0,L,A_1\sim\operatorname{Ber}(\theta)$ with
$\theta=1/2\pm b/\sqrt n$. The target is $\theta$ and the outcome
law is constant in every continuous coordinate, hence lies in every
stated smoothness ball under the explicit subfamily assumption. Product
KL is $O(b^2)$. Choose fixed $b>0$ so total variation is at most $1/2$;
the nearer-target test yields a lower bound $c/n$ for squared-error risk.
Together with (9) this proves the restricted sharp rate.
\end{proof}

For a concrete improvement, take $d_0=d_1=1$, $s_0=s_1=s_m=1$ and
$s_g=2/5$, with any positive $s_p$. Then
\[
 \gamma_0=1/3,\qquad \gamma_1=\gamma_m=1/4,\qquad \gamma_Q=2/9,
\]
so (10) holds. Separate histogram estimation of the transition density
would instead yield
\[
 \gamma_g=\frac{2/5}{2(2/5)+2}=\frac17,\qquad
 \gamma_0+\min(\gamma_m,\gamma_g)=\frac{10}{21}<\frac12,
\]
which does not satisfy the earlier sufficient condition.
The comparison concerns these proved first-order constructions; it is not
a lower bound against every estimator that explicitly estimates $g$.

\section*{Version 3: regress the doubly robust pseudo-outcome}
Version 2 regressed $\widehat m(X,L)$ to construct the first-stage
regression. Its contamination term in (12) has the rate of $\widehat m$
alone. We now replace that response by a second-stage augmented response.
This changes the proved sufficient region, while preserving the same
canonical two-stage model and the capped histogram approximation orders.
The indices $s_p,s_0,s_g,s_1,s_m$ are supplied for this minimax construction;
no adaptation to unknown indices is asserted.

Retain $t_j=\min(s_j,1)$, $D=d_0+d_1$,
$t_Q=\min(s_m,s_g,1)$ and
\[
 \gamma_0=\frac{t_0}{2t_0+d_0},\quad
 \gamma_1=\frac{t_1}{2t_1+D},\quad
 \gamma_m=\frac{t_m}{2t_m+D},\quad
 \gamma_Q=\frac{t_Q}{2t_Q+d_0}.
\]
The density assumptions used below are the canonical bounds
$0<c_p\le p\le C_p$ and $0<c_g\le g\le C_g$ on their full cubes.
They yield a lower bound for the selected design density, not merely an
integrated overlap condition. All constants below may depend on these
bounds, fixed dimensions, smoothness indices, radii and $\varepsilon$.

\begin{lemma}[Pointwise histogram error and a second-order contamination]
Construct $h_1$ and $\widehat m$ on independent blocks of $N$
trajectories, exactly as above, with side lengths tuned to $t_1,t_m$.
Then, for almost every $z=(x,l)$, uniformly in $z$ and the model,
\[
 \mathbb E|h_1(z)-e_1(z)|^2\le C N^{-2\gamma_1},\qquad
 \mathbb E|\widehat m(z)-m(z)|^2\le C N^{-2\gamma_m}.
                                                        \tag{13}
\]
On a third, independent block define
\[
 U=\widehat m(X,L)+\frac{I_1}{h_1(X,L)}
                         \{Y-\widehat m(X,L)\}.
\]
Conditionally on both fitted functions, its selected regression is
\[
 \begin{aligned}
 \mathbb E[U\mid X=x,I_0=1,h_1,\widehat m]&=Q(x)+r(x),\\
 r(x)&=\int g(l\mid x)
       \left(1-\frac{e_1(x,l)}{h_1(x,l)}\right)
                        [\widehat m(x,l)-m(x,l)]\,dl.
 \end{aligned}                                          \tag{14}
\]
Consequently
\[
 \mathbb E\|r\|_{L^2(p)}^2
                  \le C N^{-2(\gamma_1+\gamma_m)}.       \tag{15}
\]
\end{lemma}
\begin{proof}
For a histogram cell of side $h$ in dimension $D$, the probability
of selection and membership is at least
$c_p c_g\varepsilon h^D$ for $h_1$, and at least
$c_p c_g\varepsilon^2h^D$ for $\widehat m$.
The selected cell mean differs from the target regression at every
point of the cell by at most $C h^t$. Given a positive count $j$, the
sample response mean has variance at most $1/(4j)$. The binomial
reciprocal identity already proved, including the empty-cell default,
therefore bounds pointwise squared error by
$C\{h^{2t}+[(N+1)h^D]^{-1}\}$. Clipping cannot increase the error.
Substitution of the two bandwidths proves (13); importantly, no
integration over $z$ was required for this bound.

Conditional on $X,L,I_0=1$, the expectation of $I_1Y$ is $e_1m$.
Expanding $U$ and integrating over $L$ gives (14). Jensen's inequality
for the conditional density $g$, and $h_1\ge\varepsilon$, give
\[
 \|r\|_{L^2(p)}^2\le\varepsilon^{-2}
  \int (h_1-e_1)^2(\widehat m-m)^2\,d\mu.
\]
At every fixed $z$, the two squared errors are independent because their
entire training blocks are independent. Taking expectations, using
Tonelli's theorem for the nonnegative integrand, and applying the two
uniform bounds in (13) proves (15). This factorization does not assert
independence of prediction errors at different covariate points.
\end{proof}

On the independent $U$ block, among observations with $I_0=1$, regress
$U$ on $X$ by an equal-cell histogram with
$h_Q\asymp N^{-1/(2t_Q+d_0)}$, use default $1/2$ for empty cells,
and clip the fitted function to $[0,1]$. Denote it by $\widehat Q_{\rm dr}$.
The response need not lie in $[0,1]$, but it obeys
$|U|\le 1+1/\varepsilon$ deterministically.

\begin{lemma}[Risk of the augmented nested fit]
For this construction,
\[
 \mathbb E\|\widehat Q_{\rm dr}-Q\|_{L^2(p)}^2
 \le C\{N^{-2\gamma_Q}+N^{-2(\gamma_1+\gamma_m)}\}.
                                                        \tag{16}
\]
\end{lemma}
\begin{proof}
Condition on the two pilots. In an $X$-cell $A$, let
$\rho_A=\int_A e_0p$ and $\pi_A=\int_A p$.
The conditional mean of a selected response equals the selected
cell average of $Q+r$. The selected average of $Q$ differs from
$Q(x)$ throughout the cell by at most $C h_Q^{t_Q}$, by the integrated
regression modulus proved earlier. The selected average $r_A$ of $r$
satisfies
\[
 \sum_A\pi_A r_A^2
 \le\sum_A\frac{\pi_A}{\rho_A}\int_A e_0p r^2
 \le\varepsilon^{-1}\|r\|_{L^2(p)}^2.
\]
A positive-count response mean has conditional variance at most
$(1+1/\varepsilon)^2/j$. The same count calculation as in (7)
gives an integrated variance and empty-cell contribution
$C h_Q^{-d_0}/(N+1)$. Use $(a+b+c)^2\le3(a^2+b^2+c^2)$ for
sampling error, approximation and contamination; clipping reduces error
to the true $Q\in[0,1]$. Thus the conditional risk is at most
$C\{h_Q^{2t_Q}+h_Q^{-d_0}/(N+1)+\|r\|_2^2\}$.
Average over the pilots, use (15), and tune $h_Q$ to prove (16).
\end{proof}

\begin{theorem}[A strictly enlarged attained parametric region]
There is a five-block estimator, using the augmented fit above in the
longitudinal score (1), such that
\[
 \sup_{P\in\mathcal P(s_p,s_0,s_g,s_1,s_m)}
 \mathbb E_P(\widehat\psi_a-\psi_a)^2
 \le C\left\{n^{-1}+n^{-2(\gamma_0+\gamma_Q)}
                       +n^{-2(\gamma_1+\gamma_m)}\right\}.
                                                        \tag{17}
\]
Under the same constant-subfamily condition used for the lower bound
above, the squared-error minimax risk is $\Theta(n^{-1})$ whenever
\[
 \gamma_0+\gamma_Q\ge\tfrac12,\qquad
 \gamma_1+\gamma_m\ge\tfrac12.                         \tag{18}
\]
\end{theorem}
\begin{proof}
Use blocks 1, 2 and 3 to fit $h_0,h_1,\widehat m$ respectively;
block 4 to fit $\widehat Q_{\rm dr}$ from $U$; and block 5 to
average (1) with $\widehat Q=\widehat Q_{\rm dr}$ and project its
mean to $[0,1]$. All block sizes are $N=\lfloor n/5\rfloor$;
small $n$ are handled by a constant estimator.
The exact remainder (2) is valid for this new fit without alteration.
The first product in (3) factors in expectation because $h_0$ uses
only block 1 and $\widehat Q_{\rm dr}$ uses only blocks 2--4.
The second product factors because $h_1$ and $\widehat m$ use
independent blocks 2 and 3. Their mutual dependence with
$\widehat Q_{\rm dr}$ is immaterial to this sum of two bounds.
Using (11) and (16) gives the right side of (17), plus
$C n^{-2(\gamma_0+\gamma_1+\gamma_m)}$.
Since $\gamma_0>0$, this extra term is bounded by the last term in
(17). The bounded-score evaluation variance is $C/n$.
This proves (17). Conditions (18) give the upper parametric rate;
the constant Bernoulli two-point lower bound previously proved supplies
the matching lower rate in every such class.
\end{proof}

For a strict comparison take $d_0=d_1=1$,
$s_0=1/3$ and $s_1=s_m=s_g=1$, with any positive $s_p$.
Then $\gamma_0=1/5$, $\gamma_1=\gamma_m=1/4$ and
$\gamma_Q=1/3$. The new sums are $8/15$ and $1/2$, so (18)
holds, whereas version 2 additionally required
$\gamma_0+\gamma_m=9/20\ge1/2$, which fails.
This is a strict enlargement of the region established by these
constructions; it is not an impossibility theorem for other estimators
outside that region. Doubly robust augmentation is a standard methodological tool
\cite{bang}; the new contribution within this attempt is the
fully proved, index-specific improvement over its previous bound.

\section*{Unresolved minimax interactions}
The augmented nested fit removes the extra first-stage-propensity times
final-outcome-regression condition from the previous sufficient region.
The result retains the canonical unknown transition and assignment laws,
but its histogram approximation orders remain capped at one. Conditions
(18) are sufficient, not claimed necessary. Higher-order corrections,
matching slow-region lower bounds, the full five-index phase diagram and
the fixed-multistage characterization remain unresolved here. The earlier
known-assignment horizon obstruction concerns a separate growing-horizon
experiment and does not supply those missing lower bounds.

\section{Completion Estimate}
40\%

This is a heuristic estimate for the full catalogue target. Version 3 gives
a strict, constructive enlargement of the attained parametric region by
proving the second-order contamination bound and all required sample
independences. The full sharp slow-rate characterization and its
fixed-multistage extension remain missing.

\begin{thebibliography}{9}
\bibitem{agenda} C. Cinelli, A. Feller, G. Imbens, E. Kennedy, S. Magliacane and J. Zubizarreta.
\emph{Challenges in Statistics: A Dozen Challenges in Causality and Causal Inference} (2025), arXiv:2508.17099v1.
\url{https://arxiv.org/html/2508.17099v1}.
\bibitem{robins} J. Robins, L. Li, E. Tchetgen Tchetgen and A. van der Vaart.
\emph{Technical Report: Higher Order Influence Functions and Minimax Estimation of Nonlinear Functionals} (2016).
\url{https://arxiv.org/abs/1601.05820}.
\bibitem{bang} H. Bang and J. M. Robins.
\emph{Doubly Robust Estimation in Missing Data and Causal Inference Models}.
Biometrics 61 (2005), 962--973.
\end{thebibliography}
\end{document}
